% !TeX program = lualatex
\documentclass[12pt,aspectratio=169]{beamer}
\usepackage[utf8]{inputenc}
\usepackage[T1]{fontenc}
\usepackage{amsmath,amsfonts,amssymb}
\usepackage{graphicx}
\usepackage{xcolor}
\usepackage{hyperref}
\usepackage{anyfontsize}
\usepackage{lmodern}

% ====== EXTREME FONT REDUCTION ======
\renewcommand{\tiny}{\fontsize{4}{5}\selectfont}
\renewcommand{\scriptsize}{\fontsize{5}{6}\selectfont}
\renewcommand{\footnotesize}{\fontsize{6}{7}\selectfont}
\renewcommand{\small}{\fontsize{7}{8}\selectfont}
\renewcommand{\normalsize}{\fontsize{8}{9}\selectfont}
\renewcommand{\large}{\fontsize{9}{10}\selectfont}
\renewcommand{\Large}{\fontsize{10}{11}\selectfont}
\renewcommand{\LARGE}{\fontsize{11}{13}\selectfont}
\renewcommand{\huge}{\fontsize{12}{14}\selectfont}
\renewcommand{\Huge}{\fontsize{14}{17}\selectfont}

% ====== COLOR THEME ======
\definecolor{redcorrect}{RGB}{200,30,30}
\definecolor{bluecorrect}{RGB}{30,80,200}
\definecolor{greencheck}{RGB}{0,150,50}
\definecolor{lightgray}{RGB}{245,245,245}
\definecolor{darkgray}{RGB}{50,50,50}
\definecolor{softyellow}{RGB}{255,248,200}
\definecolor{steppurple}{RGB}{150,50,200}
\definecolor{deepblue}{RGB}{20,60,150}
\definecolor{darkred}{RGB}{150,20,20}
\definecolor{orangewarning}{RGB}{255,140,0}

\setbeamercolor{title}{fg=white,bg=redcorrect}
\setbeamercolor{frametitle}{fg=white,bg=redcorrect}
\setbeamercolor{block title}{fg=white,bg=bluecorrect}
\setbeamercolor{block body}{fg=darkgray,bg=lightgray}
\setbeamercolor{normal text}{fg=darkgray}
\usecolortheme{default}

% ====== CUSTOM COMMANDS ======
\newcommand{\correct}[1]{\textcolor{greencheck}{\textbf{\checkmark}} #1}
\newcommand{\keypoint}[1]{\textcolor{deepblue}{\textbf{#1}}}
\newcommand{\hardconcept}[1]{\textcolor{darkred}{\textbf{#1}}}
\newcommand{\tipbox}[1]{\fcolorbox{bluecorrect}{softyellow}{\parbox{0.88\textwidth}{\centering #1}}}
\newcommand{\stepbox}[1]{%
	\begin{center}
		\fcolorbox{darkgray}{white}{%
			\parbox{0.90\textwidth}{\vspace{0.05cm} #1 \vspace{0.05cm}}%
		}
	\end{center}
}
\newcommand{\wrongbox}[1]{\fcolorbox{redcorrect}{red!10}{\parbox{0.88\textwidth}{\centering \textcolor{redcorrect}{\textbf{\times}} #1}}}
\newcommand{\highlight}[1]{\textcolor{redcorrect}{\textbf{#1}}}
\newcommand{\bluetext}[1]{\textcolor{bluecorrect}{\textbf{#1}}}

\begin{document}
	
	% ================================================================
	% SLIDE 1: WELCOME & COURSE INTRODUCTION
	% ================================================================
	\begin{frame}
		\centering
		\vspace{1.0cm}
		{\fontsize{16}{19}\selectfont \textbf{CAMS Exam Prep (Unofficial)}}\\[0.1cm]
		{\fontsize{14}{17}\selectfont \textbf{Part B: Q\&A Styled Course}}\\[0.5cm]
		{\fontsize{10}{12}\selectfont Instructor: Shivgan Joshi}\\[0.2cm]
		{\fontsize{9}{11}\selectfont Mastering AML Frameworks and Exam-Style Questions}\\[0.5cm]
		\rule{5cm}{0.5pt}
		\vspace{0.4cm}
		\vspace{0.3cm}
		{\fontsize{8}{10}\selectfont \textcolor{darkgray}{Welcome! This course is designed to help you quickly understand and retain the most important anti-money laundering (AML) concepts tested in the CAMS exam, using question-answer and summary slides.}}
	\end{frame}
	
	% ================================================================
	% SLIDE 2: WHAT YOU WILL LEARN
	% ================================================================
	\begin{frame}
		\frametitle{\fontsize{11}{13}\selectfont What You Will Learn}
		\scriptsize
		\begin{block}{\fontsize{9}{11}\selectfont Focused CAMS Part B Topics}
			\vspace{0.05cm}
			\scriptsize
			\begin{itemize}
				\item \keypoint{FATF 40 Recommendations:} A thorough breakdown of the risk-based approach, customer due diligence, and international cooperation.
				\item \keypoint{Key Regulatory Frameworks:} Understand the US Patriot Act and the EU AML Directives, as they frequently appear in exam questions.
				\item \keypoint{Global Organizations:} Explore the roles and standards of the Egmont Group, Wolfsberg Principles, and UN Sanctions frameworks.
				\item \keypoint{Exam-Style Practice:} Reinforce your learning with realistic practice questions, explanations, and common trap identification.
			\end{itemize}
			\vspace{0.05cm}
		\end{block}
		\vspace{0.05cm}
		\tipbox{\scriptsize \textbf{Course Goal:} Provide a focused, exam-oriented review to help you confidently approach the CAMS certification.}
	\end{frame}
	
	% ================================================================
	% SLIDE 3: WHY THIS COURSE & HOW TO USE IT
	% ================================================================
	\begin{frame}
		\frametitle{\fontsize{11}{13}\selectfont Why This Course and How to Use It}
		\scriptsize
		\begin{block}{\fontsize{9}{11}\selectfont Ideal for Efficient Exam Preparation}
			\vspace{0.05cm}
			\scriptsize
			\begin{itemize}
				\item \keypoint{Designed for Success:} Created for AML analysts, compliance professionals, and certification candidates who need a targeted review of CAMS Part B.
				\item \keypoint{Last-Minute Revision:} The structured Q\&A format is perfect for quick, effective studying and understanding common exam traps.
				\item \keypoint{Build Your Knowledge:} Move step-by-step from core concepts to practical application to strengthen your foundation and readiness.
			\end{itemize}
			\vspace{0.05cm}
		\end{block}
		\vspace{0.05cm}
		\begin{center}
			\fcolorbox{bluecorrect}{softyellow}{\parbox{0.85\textwidth}{\centering \textbf{\textcolor{deepblue}{Tip:}} Use this course alongside the official CAMS study materials for a comprehensive preparation strategy.}}
		\end{center}
	\end{frame}
	
\end{document}